\documentclass[letterpaper,12pt]{article}
\usepackage[margin=0in]{geometry}
\usepackage[T1]{fontenc}
\usepackage{lmodern}
\usepackage{tikz}
\usetikzlibrary{calc,arrows.meta}
\usepackage{xcolor}
\pagestyle{empty}
\setlength{\parindent}{0pt}
\renewcommand{\familydefault}{\sfdefault}
\definecolor{quiet}{gray}{0.48}
\definecolor{rulegray}{gray}{0.76}
\definecolor{demoblue}{RGB}{29,83,139}
\definecolor{demored}{RGB}{150,54,46}

% All diagrams use the same unit in both axes. Main boards have a
% circumdiameter of four inches; small diagrams are convention examples
% or copies for saving a configuration, rather than material-fit boards.
\newcommand{\pageopen}[2]{%
  \begin{tikzpicture}[remember picture,overlay,x=1in,y=1in]
  \begin{scope}[shift={(current page.south west)}]
  \node[anchor=west,inner sep=0pt,font=\fontsize{13}{16}\selectfont]
    at (0.65,10.49) {Week 14 / Triangulations / Grades #1};
  \draw[rulegray,line width=.45pt] (0.65,10.27)--(7.85,10.27);
  \node[anchor=west,inner sep=0pt,font=\fontsize{9}{11}\selectfont]
    at (0.65,.40) {Bellingham Math Circle / Week 14 / W14-RV-draft};
  \node[anchor=east,inner sep=0pt,font=\fontsize{9}{11}\selectfont]
    at (7.85,.40) {#2};
}
\newcommand{\pageclose}{\end{scope}\end{tikzpicture}\null\newpage}
\newcommand{\textat}[3]{%
  \node[anchor=north west,inner sep=0pt,text width=7.2in,
    align=left,font=\fontsize{13}{17}\selectfont] at (#1,#2) {#3};}
\newcommand{\smalltextat}[3]{%
  \node[anchor=north west,inner sep=0pt,text width=7.2in,
    font=\fontsize{11}{14}\selectfont] at (#1,#2) {#3};}

% A at the top; A,B,C,D,E,F follow clockwise. Names stay fixed.
\newcommand{\hexagon}[3]{%
  \begin{scope}[shift={(#1,#2)}]
  \foreach \i/\lab in {0/A,1/B,2/C,3/D,4/E,5/F}{%
    \pgfmathsetmacro{\ang}{90-60*\i}
    \coordinate (h\lab) at (\ang:#3in);
    \pgfmathsetmacro{\lr}{#3+0.18}
    \node[font=\fontsize{12}{14}\selectfont,inner sep=0pt]
      at (\ang:\lr in) {\lab};
  }
  \draw[line width=.9pt] (hA)--(hB)--(hC)--(hD)--(hE)--(hF)--cycle;
}
\newcommand{\fanlines}{\draw[line width=.85pt] (hA)--(hC) (hA)--(hD) (hA)--(hE);}
\newcommand{\centrallines}{\draw[line width=.85pt] (hA)--(hC)--(hE)--cycle;}
\newcommand{\hexdots}{%
  \foreach \lab in {A,B,C,D,E,F}{\fill (h\lab) circle[radius=1.4pt];}}
\newcommand{\colorcorners}{%
  \foreach \lab in {A,B,C,D,E,F}{%
    \draw[fill=white,line width=.7pt] (h\lab) circle[radius=.12in];}}
\newcommand{\hexend}{\end{scope}}

% A pentagon convention example; its scale is equal in x and y.
\newcommand{\pentagon}[2]{%
  \begin{scope}[shift={(#1,#2)}]
  \foreach \i/\lab in {0/A,1/B,2/C,3/D,4/E}{%
    \pgfmathsetmacro{\ang}{90-72*\i}
    \coordinate (p\lab) at (\ang:.49in);
    \node[font=\fontsize{9}{11}\selectfont,inner sep=0pt]
      at (\ang:.62in) {\lab};
  }
}
\newcommand{\peelrow}[1]{%
  \node[anchor=west,inner sep=0pt,font=\fontsize{12}{15}\selectfont]
    at (.75,#1) {Peeled corners:};
  \draw[rulegray,line width=.5pt] (2.08,#1-.08)--(4.50,#1-.08);
  \node[anchor=west,inner sep=0pt,font=\fontsize{12}{15}\selectfont]
    at (4.83,#1) {Last triangle:};
  \draw[rulegray,line width=.5pt] (6.04,#1-.08)--(7.73,#1-.08);
}
\newcommand{\savehex}[2]{%
  \hexagon{#1}{#2}{.81}
  \hexdots
  \hexend
  \node[anchor=west,inner sep=0pt,font=\fontsize{11}{13}\selectfont]
    at (#1-.95,#2-1.13) {Turn:};
  \draw[rulegray,line width=.5pt] (#1-.48,#2-1.18)--(#1+.96,#2-1.18);
}

\begin{document}

\pageopen{K--1}{1}
\textat{.65}{10.03}{An ear is a triangle with two sides on the outside edge.
Peel it by erasing those sides and their meeting corner. Keep its third side.
Trace a drawing, then lift the tracing off this page before peeling.}

% Worked convention: input, removal, output; a different polygon from
% the investigation boards, with the removed corner explicitly identified.
\pentagon{1.85}{8.53}
  \draw[line width=.8pt] (pA)--(pB)--(pC)--(pD)--(pE)--cycle;
  \draw[line width=.7pt] (pA)--(pC) (pA)--(pD);
\end{scope}
\pentagon{4.25}{8.53}
  \fill[demoblue!15] (pA)--(pB)--(pC)--cycle;
  \draw[line width=.8pt] (pA)--(pB)--(pC)--(pD)--(pE)--cycle;
  \draw[line width=.7pt] (pA)--(pC) (pA)--(pD);
  \draw[demored,line width=1.1pt] ($(pB)+(-.09,-.09)$)--($(pB)+(.09,.09)$)
    ($(pB)+(-.09,.09)$)--($(pB)+(.09,-.09)$);
  \draw[demored,line width=1.1pt] (pA)--(pB)--(pC);
\end{scope}
\begin{scope}[shift={(6.65,8.53)}]
  \foreach \i/\lab in {0/A,2/C,3/D,4/E}{%
    \pgfmathsetmacro{\ang}{90-72*\i}
    \coordinate (p\lab) at (\ang:.49in);
    \node[font=\fontsize{9}{11}\selectfont,inner sep=0pt]
      at (\ang:.62in) {\lab};
  }
  \draw[line width=.8pt] (pA)--(pC)--(pD)--(pE)--cycle;
  \draw[line width=.7pt] (pA)--(pD);
\end{scope}
\foreach \x/\lab in {1.85/Before,4.25/Peel B,6.65/After}{%
  \node[font=\fontsize{11}{14}\selectfont,inner sep=0pt] at (\x,7.83) {\lab};}
\draw[-{Stealth[length=4pt]},quiet,line width=.6pt] (2.64,8.53)--(3.40,8.53);
\draw[-{Stealth[length=4pt]},quiet,line width=.6pt] (5.04,8.53)--(5.80,8.53);

\textat{.65}{7.43}{\textbf{Problem 1:} Peel ears from each large drawing on these two pages until one
triangle remains, writing the peeled corner letters in order. Can your choices
leave different last triangles?}
\hexagon{4.25}{4.54}{2}
  \fanlines\hexdots
\hexend
\peelrow{1.90}
\peelrow{1.29}
\pageclose

\pageopen{K--1}{2}
\textat{.65}{10.03}{Extra workspace for Problem 1}
\hexagon{4.25}{7.37}{2}
  \centrallines\hexdots
\hexend
\peelrow{4.53}
\peelrow{3.84}
\peelrow{3.15}
\peelrow{2.46}
\peelrow{1.77}
\pageclose

\pageopen{2--3}{3}
\textat{.65}{10.03}{Use red, blue and yellow. Give each corner one color.}
\textat{.65}{9.43}{\textbf{Problem 2:} Color the corners of each large drawing on these two pages so
every triangle has all three colors. Circle as few corners as possible so every
triangle touches a circled corner.}
\hexagon{4.25}{6.27}{2}
  \fanlines\colorcorners
\hexend
\smalltextat{.65}{3.58}{Extra workspace}
\hexagon{2.55}{2.03}{1.08}
  \fanlines\colorcorners
\hexend
\hexagon{5.95}{2.03}{1.08}
  \fanlines\colorcorners
\hexend
\pageclose

\pageopen{2--3}{4}
\textat{.65}{10.03}{Extra workspace for Problem 2}
\hexagon{4.25}{7.17}{2}
  \centrallines\colorcorners
\hexend
\hexagon{2.55}{2.85}{1.08}
  \centrallines\colorcorners
\hexend
\hexagon{5.95}{2.85}{1.08}
  \centrallines\colorcorners
\hexend
\pageclose

\pageopen{4--5}{5}
\textat{.65}{10.03}{Join corners with straight lines to divide the hexagon into
triangles. Lines may meet only at corners. Keep the corner letters fixed;
different sets of lines count as different drawings.}
\textat{.65}{9.23}{Trace the inside lines and the center mark. Turn the tracing
about that mark. It matches when every turned line lands on an original line.}

% Worked rotation convention: one diagonal, deliberately not a complete
% hexagon triangulation. Fixed corner labels show A -> D and C -> F.
\hexagon{1.85}{7.97}{.47}
  \draw[demoblue,line width=1.1pt] (hA)--(hC);
  \draw[quiet,line width=.5pt] (-.035,0)--(.035,0) (0,-.035)--(0,.035);
\hexend
\hexagon{4.25}{7.97}{.47}
  \draw[demoblue!40,line width=.8pt,dashed] (hA)--(hC);
  \draw[demored,line width=1.1pt] (hD)--(hF);
  \draw[-{Stealth[length=3.5pt]},quiet,line width=.7pt]
    (90:.20in) arc[start angle=90,end angle=-90,radius=.20in];
  \draw[quiet,line width=.5pt] (-.035,0)--(.035,0) (0,-.035)--(0,.035);
\hexend
\hexagon{6.65}{7.97}{.47}
  \draw[demored,line width=1.1pt] (hD)--(hF);
  \draw[quiet,line width=.5pt] (-.035,0)--(.035,0) (0,-.035)--(0,.035);
\hexend
\foreach \x/\lab in {1.85/Before,4.25/Half-turn,6.65/After}{%
  \node[font=\fontsize{11}{14}\selectfont,inner sep=0pt] at (\x,7.03) {\lab};}
\draw[-{Stealth[length=4pt]},quiet,line width=.6pt] (2.61,7.97)--(3.49,7.97);
\draw[-{Stealth[length=4pt]},quiet,line width=.6pt] (5.01,7.97)--(5.89,7.97);

\textat{.65}{6.65}{\textbf{Problem 3:} Find all triangulations that match after
a half-turn, a one-third-turn or a one-sixth-turn. Save each with its turn on the
next page. Explain why each list is complete.}
\hexagon{4.25}{3.57}{2}
  \hexdots
  \draw[quiet,line width=.6pt] (-.06,0)--(.06,0) (0,-.06)--(0,.06);
\hexend
\pageclose

\pageopen{4--5}{6}
\textat{.65}{10.03}{Extra workspace for Problem 3}
\foreach \y in {8.62,6.17,3.72}{%
  \foreach \x in {1.75,4.25,6.75}{\savehex{\x}{\y}}}
\smalltextat{.65}{2.15}{Explanation:}
\draw[rulegray,line width=.45pt] (.65,1.58)--(7.85,1.58);
\draw[rulegray,line width=.45pt] (.65,1.12)--(7.85,1.12);
\end{scope}\end{tikzpicture}\null

\end{document}
